\documentclass[12pt,a4paper]{article}
\usepackage[utf8]{inputenc}
\usepackage[T1]{fontenc}
\usepackage[polish]{babel}
\usepackage{geometry}
\geometry{margin=2.5cm}

\begin{document}

\section*{Raport z modelowania: Prognozowanie popytu i analiza szeregów czasowych w systemie wieloagentowym}

Zbudowanie skutecznego modelu predykcyjnego dla środowiska e-commerce wymaga przekształcenia wysoce zaszumionych logów transakcyjnych w ustrukturyzowany problem uczenia maszynowego. W pierwszej kolejności surowe dane sprzedażowe poddawane są procesowi agregacji do interwałów dziennych, co pozwala na wstępną redukcję mikrozmian i stabilizację sygnału. Zanim jednak algorytmy rozpoczną pracę, system przeprowadza rygorystyczny audyt uczalności (learnability) asortymentu. Na podstawie liczby dni sprzedaży, odchylenia standardowego oraz udziału dni z zerową sprzedażą, produkty kategoryzowane są na trzy grupy: minimalną, praktyczną oraz rygorystyczną (Strict), zrzeszającą bestsellery o najsilniejszym sygnale. Kategoryzacja ta pozwala na alokację zasobów obliczeniowych tylko tam, gdzie sztuczna inteligencja faktycznie pokonuje prostą średnią ruchomą. Następnie przeprowadzana jest inżynieria cech. Ponieważ modele opierające się na drzewach decyzyjnych nie potrafią natywnie interpretować czasu, data ulega dekompozycji, a dodatkowo budowane są okna opóźnień oraz statystyki kroczące. Kolejnym krytycznym krokiem jest świadomość braków magazynowych (stockout awareness). System inteligentnie odróżnia brak sprzedaży wynikający z pustego magazynu od braku zainteresowania klientów, co jest niezbędne do modelowania faktycznego, a nie stłumionego popytu. Ostatnim etapem przygotowań jest nałożenie limitów na skrajne odchylenia (outlier capping) przy użyciu mnożnika 3.0 dla rozstępu międzykwartylowego (IQR). Co niezwykle istotne dla zachowania rygoru eksperymentu, widełki te wyliczane są wyłącznie na zbiorze treningowym, co całkowicie eliminuje zjawisko wycieku danych z przyszłości (data leakage).

Niezwykle istotnym aspektem przygotowania danych jest ich właściwy podział na zbiory uczące, walidacyjne i testowe, realizowany w bezpiecznej proporcji 70/15/15. W przypadku analizy szeregów czasowych kategorycznie zabrania się stosowania standardowej, losowej walidacji krzyżowej. Aby uniknąć krytycznego błędu, system stosuje rygorystyczny podział chronologiczny, trenując model wyłącznie na przeszłości. Dodatkową warstwą ochronną na tym etapie jest analiza przesunięcia rozkładów (distribution shift). System automatycznie diagnozuje zgodność współczynnika zmienności pomiędzy zbiorem uczącym a testowym. Jeśli dystrybucje różnią się o ponad 30 procent, algorytm generuje ostrzeżenie o drastycznej zmianie otoczenia rynkowego względem historii. W tym kontekście należy również zrozumieć fundamentalną rolę wariancji. Choć w klasycznej statystyce wariancja często utożsamiana jest z błędem, w uczeniu maszynowym stanowi główny nośnik informacji. Fluktuacje sprzedaży tworzą złożoną strukturę, w której algorytmy poszukują wzorców matematycznych. Szeregi, w których zbiór testowy nie wykazuje minimum wymaganej wariancji, są automatycznie odrzucane z procesu ewaluacji, co zapobiega testowaniu algorytmów na ,,płaskich'' danych.

W silniku predykcyjnym zaimplementowano trzy diametralnie różne architektury modeli, z których każda charakteryzuje się unikalnym podejściem do analizy danych. Pierwszą z nich jest Extreme Gradient Boosting, w skrócie XGBoost. Stanowi on klasyczne i niezwykle solidne rozwiązanie do analizy danych tabelarycznych, które optymalizuje sekwencyjne budowanie płytkich drzew decyzyjnych w celu minimalizacji błędów generowanych przez poprzedników. Podczas fitowania modeli drzewiastych zastosowano genialną mechanikę czyszczenia celu: filtrując dni, w których magazyn był pusty, modele uczą się wyłącznie w oparciu o pełną dostępność towaru. W efekcie algorytmy odtwarzają surowy i nieskrępowany popyt rynkowy, a nie sztucznie stłumioną sprzedaż. Ponadto, zamiast klasycznego błędu średniokwadratowego (RMSE), optymalizują one uodpornioną na skrajne wartości odstające funkcję błędu Hubera (oraz pseudo-Hubera), która w przypadku ogromnych odchyleń ulega linearyzacji. Zapobiega to gwałtownemu wariowaniu modelu podczas nienaturalnych, ogromnych skoków sprzedaży.

Drugą architekturą jest Light Gradient Boosting Machine, opracowany przez firmę Microsoft jako lżejsza i znacznie szybsza alternatywa dla klasycznego XGBoosta. Jego kluczowa różnica polega na asymetrycznym budowaniu drzew. Zamiast rozwijać się poziomami, LightGBM rośnie w głąb liści, zawsze wybierając ten węzeł, który gwarantuje największą redukcję błędu. Ta cecha sprawia, że algorytm ten jest fenomenalnie wydajny pamięciowo i obliczeniowo, co nabiera krytycznego znaczenia przy skalowaniu systemu na tysiące produktów o rzadszej sprzedaży. Wymaga on jednak znacznie agresywniejszej regularyzacji, aby zapobiec zapamiętywaniu szumu w mniejszych zbiorach danych.

Trzecim podejściem jest Long Short-Term Memory, czyli architektura LSTM należąca do rodziny Głębokich Sieci Neuronowych. W przeciwieństwie do modeli drzewiastych, sieci LSTM zostały stworzone specjalnie do natywnego analizowania sekwencji. Wymagają one jednak zupełnie innej transformacji danych. Zamiast wyciągać dyskretne cechy czasu, dane formuje się w trójwymiarowe, przesuwające się okna czasowe. Z uwagi na wrażliwość sieci na rozstrzał wartości, wprowadzono rygorystyczną normalizację danych. Trening takiego modelu jest wspierany przez proces zautomatyzowanego strojenia hiperparametrów za pomocą biblioteki Optuna. Algorytm samoczynnie poszukuje optymalnej szerokości okna czasowego (np. 7, 14 lub 21 dni), liczby neuronów, współczynnika dropout oraz idealnego tempa uczenia. Zastosowano również redukcję tempa uczenia na wypadek stagnacji ewaluacyjnej (ReduceLROnPlateau), co gwarantuje płynne i bezbłędne zbiegnięcie do minimum kosztu.

Zamiast arbitralnie narzucać jeden algorytm dla całego asortymentu, system wykorzystuje strategię turniejową na poziomie pojedynczego produktu. Modele walczą między sobą o miano mistrza na odłożonym zbiorze walidacyjnym. Dla przypadków wyjątkowo trudnych, charakteryzujących się rzadką lub mocno rwaną sprzedażą, system używa pre-kalkulowanego słownika strategii zapasowych (fallback cache). Obejmuje on mapowanie korelacji rówieśniczej (peer-correlation), gdzie popyt modelowany jest poprzez regresję Ridge względem silnie powiązanego produktu. Wdrożono także tradycyjne metody Crostona oraz TSB (Teunter-Syntetos-Babai), stworzone specjalnie z myślą o wolno rotujących zapasach, które służą jako koło ratunkowe, gdy AI zawodzi. Aby zintegrować całość ze środowiskiem produkcyjnym, system domyka rygorystyczny mechanizm bezpieczników (Guardrails). Jeśli wyliczony na zbiorze testowym błąd R-kwadrat okazuje się ujemny, a sztuczna inteligencja radzi sobie gorzej niż płaska średnia, system arbitralnie podmienia model na strategię predykcji opartą na poprzednim dniu (lag1 posthoc). Wdrożono również kalibrację dla wyjątkowo trudnych przypadków (hard cases): jeśli model globalnie niedoszacowuje lub przeszacowuje wolumen o ponad 12 procent na ewaluacji (bias ilościowy), nakładany jest mnożnik korygujący. Całość dopełnia twardy limit (anti-spike cap) odcinający piki powyżej 90. percentyla z historii, chroniący dział zaopatrzenia przed absurdalnymi wolumenami i w konsekwencji przed groźnym przeinwestowaniem. Mechanizmy te sprawiają, że moduł jest bezpiecznym, samooptymalizującym się systemem gotowym na wdrożenie produkcyjne.

Przeprowadzone symulacje obejmowały horyzonty jednego, dwóch oraz trzech miesięcy, dostarczając niezwykle cennych obserwacji. W perspektywie jedno- i dwumiesięcznej modele bazujące na drzewach decyzyjnych wykazywały się spektakularną precyzją w odtwarzaniu rytmów tygodniowych. Ich predykcje idealnie wpasowywały się w historyczne piki, tworząc charakterystyczny, grzebieniasty kształt wykresu. Z kolei sieci LSTM miały tendencję do generowania znacznie gładszych, bardziej opływowych krzywych, które z jednej strony ignorowały mikroskoki, ale z drugiej doskonale podążały za szerszym makrotrendem. Sytuacja ulegała jednak dramatycznej zmianie w perspektywie trzymiesięcznej, obnażając największą słabość modeli typu Gradient Boosting. Algorytmy te są matematycznie niezdolne do ekstrapolacji trendów liniowych wychodzących poza historyczne maksimum. Jeśli dany produkt znajdował się w silnym trendzie wzrostowym, XGBoost ,,odbijał się od sufitu'' swoich historycznych doświadczeń, podczas gdy LSTM potrafił kontynuować rysowanie trajektorii wznoszącej.

Podsumowując wyniki dla dwudziestu najlepiej sprzedających się produktów, eksperyment dowiódł, że popyt dla tej grupy można przewidzieć z bardzo wysoką precyzją, osiągając błąd procentowy na poziomie od dziesięciu do piętnastu procent w horyzoncie miesięcznym. Ten spektakularny sukces analityczny wynika bezpośrednio z wysokiego stosunku sygnału do szumu. Przy wolumenach rzędu kilkudziesięciu lub kilkuset sztuk dziennie, zgodnie z prawem wielkich liczb, indywidualne losowe decyzje konsumentów znoszą się nawzajem. Szum statystyczny ulega wygaszeniu, a na wierzch wypływają silne, łatwe do zamodelowania trendy i sezonowości. Dodatkowo, bestsellery charakteryzują się ciągłością sprzedaży, co eliminuje niezwykle trudny problem popytu rwanego. Wdrożona architektura wieloagentowa skutecznie automatyzuje i optymalizuje cały proces, stanowiąc niezawodny fundament decyzyjny dla polityki zakupowej i cenowej e-commerce.

\end{document}
